The four-pronged zero germ of Theorem~\ref{thm:q4-hecke-zero-network}
suggests a knot-theoretic question, but it is not itself a knot.  The germ lives
in a real surface, and its intersection with a sufficiently small circle consists
of four points.  A knot appears only after passing to the three-dimensional unit
tangent bundle of the global Hecke orbifold.  This section makes that passage and
keeps the local and global constructions separate.

\subsection{The sign cover as an orbifold}

Write
\[
 \mathcal O_4=\Gamma_4\backslash\HH,
 \qquad
 R=JT,
 \qquad
 H_4=\ker\chi.
\]
Thus
\begin{equation}
 \Gamma_4=\langle J,R\mid J^2=R^4=1\rangle,
 \qquad
 \chi(J)=\chi(R)=-1,
 \qquad
 T=JR\in H_4.
 \label{eq:q4-sign-kernel-presentation}
\end{equation}
The quotient
\(\mathcal O_4^\chi=H_4\backslash\HH\) is the orbifold on which the
assignment \(a_4\) of Proposition~\ref{prop:q4-sign-descent} becomes scalar.

\begin{proposition}[Signature of the sign cover]
\label{prop:q4-sign-cover-signature}
The index-two orbifold cover
\[
 \mathcal O_4^\chi\longrightarrow\mathcal O_4
\]
has signature
\[
 (0;2;\infty,\infty).
\]
More precisely, the order-two point of \(\mathcal O_4\) has one regular
preimage, the order-four point has one preimage of order two, and the cusp has
two unramified preimages.
\end{proposition}

\begin{proof}
The relevant stabilizer intersections follow immediately from
\eqref{eq:q4-sign-kernel-presentation}:
\[
 H_4\cap\langle J\rangle=1,
 \qquad
 H_4\cap\langle R\rangle=\langle R^2\rangle,
 \qquad
 \langle T\rangle\subset H_4.
\]
Since \(J\) and \(R\) lie outside \(H_4\), the first two orbifold points each
have one preimage.  Their new stabilizer orders are respectively \(1\) and
\(2\).  Since the entire parabolic stabilizer lies in \(H_4\), the number of
cusp preimages is
\([\Gamma_4:H_4\langle T\rangle]=2\), and both peripheral covering degrees
are one.

It remains only to determine the genus.  The orbifold Euler characteristic of
the base is
\[
 \chi_{\mathrm{orb}}(\mathcal O_4)
 =2-1-\left(1-\frac12\right)-\left(1-\frac14\right)
 =-\frac14.
\]
It is multiplied by two under the cover.  A genus-\(g\) orbifold with one
order-two point and two cusps has Euler characteristic
\[
 2-2g-2-\left(1-\frac12\right)=-2g-\frac12.
\]
Equating this with \(-1/2\) gives \(g=0\).
\end{proof}

\begin{proposition}[Global scalar-cover uniformization]
\label{prop:q4-sign-cover-cylinder}
Let \(\widehat{\mathcal O}_4^\chi\) be the coarse AEG surface obtained by
adjoining the finite-distance point over the order-four orbit to the regular
locus of \(H_4\backslash\HH\), while leaving its two cusps open.  Then \(W\)
descends to a biholomorphism
\begin{equation}
 W:\widehat{\mathcal O}_4^\chi\longrightarrow\C^\times.
 \label{eq:q4-sign-cover-cstar}
\end{equation}
Under this identification,
\[
 g_4=\left|\frac{\dd W}{W}\right|^2
\]
is the complete flat cylindrical metric and the scalar zero set is the single
circle
\[
 Z(a_4)/H_4=\{|W|=1\}.
\]
\end{proposition}

\begin{proof}
Proposition~\ref{prop:q4-sign-cover-signature} identifies the compactified
coarse surface with \(\mathbb P^1\), with two punctures coming from the cusps.
The \(H_4\)-invariance of \(F\) makes \(W\) single-valued on this surface.
It has neither a zero nor a pole in the interior.  At the two cusps the
expansion in the proof of Theorem~\ref{thm:q4-hecke-zero-network} is
\[
 W=Cq^{\pm1}(1+O(q)).
\]
The two signs are opposite: an element of
\(\Gamma_4\setminus H_4\) interchanges the two cusp lifts and sends
\(W\) to \(W^{-1}\).  Hence the meromorphic extension of \(W\) to the
compact coarse sphere has one simple zero and one simple pole, so it has degree
one.

At the remaining order-two orbifold point of the cover, an upstairs coordinate
\(z\) satisfies \(W=1+cz^2+O(z^4)\), with \(c\ne0\).  Thus in the coarse
coordinate \(z^2\), \(W\) is unramified.  It is also unramified at the regular
preimage of the old order-two point and everywhere else.  This proves
\eqref{eq:q4-sign-cover-cstar}.  The metric and zero-circle statements now
follow directly, and completeness is the divergence of
\(\int |\dd W|/|W|\) at \(0\) and \(\infty\).
\end{proof}

\begin{remark}[The zero circle is not a hyperbolic periodic orbit]
\label{rem:q4-zero-circle-not-geodesic-knot}
The circle in Proposition~\ref{prop:q4-sign-cover-cylinder} is a zero level and
a geodesic for the flat AEG metric.  No claim is made that it is a closed
geodesic for the original hyperbolic orbifold metric.  The periodic-orbit knots
below instead come from unit-tangent lifts of hyperbolic axes.  In particular,
Proposition~\ref{prop:q4-sign-cover-cylinder} forgets the residual order-two
orbifold isotropy in the coarse coordinate, whereas $T^1\mathcal O_4^\chi$
below retains that isotropy.  The flat unit tangent bundle of the coarse
cylinder is not being identified with the Hecke unit tangent bundle.
\end{remark}

\subsection{Hyperbolic compactification and the cusp link}

Let
\[
 M_4=T^1\mathcal O_4,
 \qquad
 M_4^\chi=T^1\mathcal O_4^\chi.
\]
The unit tangent bundle of an orientable hyperbolic orbifold is a genuine
oriented three-manifold: the rotations at an orbifold point act freely on unit
tangent vectors.

\begin{definition}[Hyperbolic cusp compactification]
For a cusp of an oriented hyperbolic two-orbifold, the \emph{hyperbolic
compactification} of its unit tangent bundle is the Dehn filling in which the
circle of vectors tangent to a horocycle bounds a meridian disc.  The core of
the attached solid torus is denoted by \(K_\infty\).  This convention is part of
the definition; changing the filling slope can change the ambient lens space.
\end{definition}

\begin{theorem}[The \(q=4\) cusp knot and its sign-cover link]
\label{thm:q4-sign-cover-torus-link}
With the preceding compactification convention:
\begin{enumerate}
\item the compactification of \(M_4\) is
\begin{equation}
 \overline M_4\cong L(2,1)\cong\mathbb {RP}^3,
 \qquad
 M_4\cong\mathbb {RP}^3\setminus K_\infty,
 \label{eq:q4-rp3-cusp-complement}
\end{equation}
and \(K_\infty\) is null-homologous in \(\mathbb {RP}^3\);
\item the covering \(M_4^\chi\to M_4\) extends over the cusp fillings to the
universal double cover
\[
 S^3\longrightarrow\mathbb {RP}^3;
\]
\item the full inverse image of \(K_\infty\) is the two-component torus link
\(T(2,4)\), and hence
\begin{equation}
 M_4^\chi\cong S^3\setminus T(2,4).
 \label{eq:q4-sign-cover-link-complement}
\end{equation}
The two components have absolute mutual linking number \(2\).
\end{enumerate}
\end{theorem}

\begin{proof}
For the triangle orbifold of signature \((p,q,\infty)\), the standard unit-tangent
bundle calculation gives
\begin{equation}
 \overline{T^1\mathcal O_{p,q}}
 \cong L(pq-p-q,p-1).
 \label{eq:dehornoy-lens-compactification}
\end{equation}
In the same calculation the added cusp fiber lies on the median Heegaard torus.
If \((c_P,c_Q)\) is the basis used there, the two meridians and the cusp fiber
have classes
\begin{equation}
 m_P=(p-1,-1),
 \qquad
 m_Q=(-1,q-1),
 \qquad
 a_\infty=(1,1).
 \label{eq:dehornoy-median-slopes}
\end{equation}
These statements are the lens-space and slope computation of
\cite[Lemmas~5.2 and~5.3]{Dehornoy2015GeodesicFlow}.  Setting \(p=2\) and \(q=4\)
gives
\(L(2,1)\cong\mathbb {RP}^3\), and deleting the added core recovers \(M_4\).
Moreover, for these values,
\[
 a_\infty=2m_P+m_Q.
\]
Both meridians vanish after the two solid tori are attached, so
\([K_\infty]=0\) in \(H_1(\mathbb {RP}^3;\Z)\).

By Proposition~\ref{prop:q4-sign-cover-signature}, each cusp has two lifts and
each peripheral covering degree is one.  More explicitly, both the parabolic
class \(T\) and the central unit-tangent fiber class lie in the covering
subgroup.  Therefore every peripheral torus maps homeomorphically to the base
peripheral torus and the chosen meridian slope lifts unchanged.  The double
cover consequently extends, without branching, over the two solid tori used
for hyperbolic filling.  The extended cover is connected, so it is the universal double cover
\(S^3\to\mathbb {RP}^3\).

It remains to identify the lifted link; this is where the non-coprime pair
\((2,4)\) requires an explicit slope check.  For \(p=2,q=4\), equations
\eqref{eq:dehornoy-median-slopes} read
\[
 m_P=(1,-1),
 \qquad m_Q=(-1,3),
 \qquad a_\infty=(1,1).
\]
The compactified first homology is
$\Z^2/\langle m_P,m_Q\rangle$, whose order is
\(|\det(m_P,m_Q)|=2\).  The homomorphism
\((x,y)\mapsto x+y\pmod 2\) kills both meridians, so it is the restriction to
the median torus of the quotient map to
\(H_1(\mathbb {RP}^3;\Z)\cong\Z/2\Z\).  Therefore the universal cover of the
median torus corresponds to the index-two lattice
\[
 \Lambda=\{(x,y)\in\Z^2:x+y\equiv0\pmod2\}.
\]
The pair \((a_\infty,m_P)\) is a basis of \(\Lambda\), and in that basis
\[
 m_P=(0,1),
 \qquad
 m_Q=(1,-2),
 \qquad
 a_\infty=(1,0).
\]
Thus the lifted Heegaard splitting is the standard genus-one splitting of \(S^3\).
Moreover \(a_\infty=m_Q+2m_P\).  Since \(a_\infty\) is trivial in the deck
quotient, its inverse image has two components, each mapping homeomorphically
to \(K_\infty\).  On the lifted Heegaard torus these are two parallel copies of
the primitive slope \(2m_P+m_Q\).  They form \(T(4,2)\), which is isotopic to
\(T(2,4)\).  The usual torus framing gives absolute mutual linking number \(2\);
its sign depends on the orientation convention.
Removing the link proves \eqref{eq:q4-sign-cover-link-complement}.
\end{proof}

\begin{warning}[What has and has not produced the link]
\label{warn:q4-local-global-link}
Theorem~\ref{thm:q4-sign-cover-torus-link} is global.  It uses the
\((2,4,\infty)\) orbifold, its unit tangent bundle, the character cover, and a
specified cusp filling.  It does not identify the small link of the real
four-pronged surface germ with \(T(2,4)\).  Nor can one obtain the conclusion by
substituting \(p=2,q=4\) into an \(S^3\) torus-\emph{knot} theorem that assumes
\(p\) and \(q\) coprime.  In particular, the central-extension presentation
\(\langle x,y\mid x^2=y^4\rangle\) has abelianization
\(\Z\oplus\Z/2\Z\), whereas the complement of a two-component link in \(S^3\)
has first homology \(\Z^2\).  The passage to the index-two cover in
Theorem~\ref{thm:q4-sign-cover-torus-link} is essential.
\end{warning}

\subsection{Primitive Hecke classes and geodesic knots}

\begin{definition}[Oriented Hecke geodesic knot]
Let \(\gamma\in\Gamma_4\) be primitive and hyperbolic.  Its invariant axis
\(A_\gamma\subset\HH\) is oriented from the repelling to the attracting fixed
point.  The corresponding closed geodesic in \(\mathcal O_4\) is denoted
\(C_\gamma\).  Its unit-tangent lift is
\begin{equation}
 K_\gamma
 =\{(x,\dot C_\gamma(x)):x\in C_\gamma\}
 \subset M_4.
 \label{eq:q4-unit-tangent-knot}
\end{equation}
\end{definition}

\begin{proposition}[Arithmetic class to periodic-orbit knot]
\label{prop:q4-hyperbolic-class-knot}
Primitive hyperbolic conjugacy classes in \(\Gamma_4\) are in bijection with
primitive oriented periodic orbits of the geodesic flow on \(M_4\).  Under this
correspondence \(K_\gamma\) is an embedded oriented circle.  Replacing
\(\gamma\) by \(\gamma^{-1}\) reverses the orientation of the base geodesic and
in general gives the oppositely directed periodic orbit rather than the same
oriented knot.

If
\[
 \widetilde\gamma=
 \begin{pmatrix}a&b\\c&d\end{pmatrix}
 \in\operatorname{SL}_2(\Q(\sqrt2))
\]
is a lift, then its fixed points satisfy
\begin{equation}
 cx^2+(d-a)x-b=0,
 \qquad
 \Delta_\gamma=(\operatorname{tr}\widetilde\gamma)^2-4,
 \label{eq:q4-axis-quadratic-data}
\end{equation}
and the geodesic length is
\begin{equation}
 \ell(K_\gamma)
 =2\operatorname{arcosh}
 \frac{|\operatorname{tr}\widetilde\gamma|}{2}.
 \label{eq:q4-geodesic-length-trace}
\end{equation}
Thus the knot is canonically labelled by quadratic Hecke data, although its
ambient isotopy type need not be determined by the trace or discriminant alone.
\end{proposition}

\begin{proof}
A hyperbolic element acts by translation along its oriented axis.  Taking the
axis modulo the cyclic subgroup it generates gives a closed geodesic; conjugation
only changes its lift in \(\HH\).  Conversely, lifting an oriented closed geodesic
to \(\HH\) gives an axis and its primitive deck translation.  Primitivity removes
multiple traversal.  A primitive periodic orbit of a nonsingular flow is an
embedded circle, even when its projection to the base orbifold self-intersects.
The fixed-point equation follows by solving
\(x=(ax+b)/(cx+d)\), and the trace--length formula is the standard eigenvalue
calculation for a hyperbolic element of \(\operatorname{PSL}_2(\R)\).
\end{proof}

The sign character has an immediate three-dimensional meaning for these knots.

\begin{proposition}[Sign-controlled lifting]
\label{prop:q4-sign-controlled-knot-lift}
Let \(\pi:M_4^\chi\to M_4\) be the sign cover and let \(K_\gamma\) be as above.
Then
\[
 \pi^{-1}(K_\gamma)
 \quad\text{has}\quad
 \begin{cases}
 2\text{ components, each of degree }1,&\chi(\gamma)=+1,\\
 1\text{ component of degree }2,&\chi(\gamma)=-1.
 \end{cases}
\]
In the second case the lifted periodic orbit is represented by
\(\gamma^2\in H_4\).
\end{proposition}

\begin{proof}
The monodromy of the double cover around the periodic orbit is the image of its
projected deck transformation in
\(\Gamma_4/H_4\cong\{\pm1\}\), namely \(\chi(\gamma)\).  The two alternatives
are the elementary lifting alternatives for a connected double cover.  If the
monodromy is nontrivial, traversing the base orbit twice closes the lift, and the
corresponding element is \(\gamma^2\).
\end{proof}

\subsection{The classical cutting code and a linking formula}

We next record the precise part of the arithmetic--knot relation that is already
available from the classical Hecke tessellation.  Choose the standard adapted
\((2,4,\infty)\) tessellation and the orientations used in
\cite{Dehornoy2015GeodesicFlow}.  A periodic geodesic has a reduced cyclic
cutting code
\begin{equation}
 \bigl(u v^{j_1}u v^{j_2}\cdots u v^{j_m}\bigr)^{\Z},
 \qquad j_k\in\{1,2,3\},
 \label{eq:q4-reduced-cutting-code}
\end{equation}
unique up to cyclic permutation.
Because \(K_\infty\) is null-homologous by
Theorem~\ref{thm:q4-sign-cover-torus-link}, its linking number with any disjoint
oriented knot is defined using an integral Seifert surface and is integer-valued.

\begin{theorem}[The \(q=4\) cusp-linking formula]
\label{thm:q4-cusp-linking-code}
For a periodic orbit \(K_\gamma\) with reduced code
\eqref{eq:q4-reduced-cutting-code}, the integer linking number in
\(\overline M_4\cong\mathbb {RP}^3\) is
\begin{equation}
 \operatorname{Lk}_{\overline M_4}(K_\gamma,K_\infty)
 =\sum_{k=1}^m(j_k-2)
 =\#\{k:j_k=3\}-\#\{k:j_k=1\}.
 \label{eq:q4-cusp-linking-specialized}
\end{equation}
The sign is tied to the displayed orientation convention; reversing one component
reverses the sign.
\end{theorem}

\begin{proof}
For a \((p,q,\infty)\) orbifold, define the wheel turn of a reduced code
\(u^{i_1}v^{j_1}\cdots u^{i_m}v^{j_m}\) by
\[
 \Theta_{\mathrm{wheel}}
 =\sum_{k=1}^m
 \left(\frac{i_k-p/2}{p}+\frac{j_k-q/2}{q}\right).
\]
The classical cusp-linking theorem gives
\[
 \operatorname{Lk}(K_\gamma,K_\infty)
 =\frac{pq}{pq-p-q}\,\Theta_{\mathrm{wheel}}
 \qquad\text{\cite[Proposition~5.7]{Dehornoy2015GeodesicFlow}}.
\]
For \(p=2\), every \(i_k\) equals \(1\), so its contribution vanishes.  For
\(q=4\), the remaining contribution is
\(\Theta_{\mathrm{wheel}}=\frac14\sum_k(j_k-2)\), while
\(pq/(pq-p-q)=4\).  Substitution proves
\eqref{eq:q4-cusp-linking-specialized}.
\end{proof}

\begin{corollary}[Parity of linking and splitting in the sign cover]
\label{cor:q4-linking-sign-parity}
For the geometric cutting code
\eqref{eq:q4-reduced-cutting-code}, the standard turn-to-deck-transformation
calibration gives
\begin{equation}
 \chi(\gamma)
 =(-1)^{m+\sum_k j_k}
 =(-1)^{m+\operatorname{Lk}(K_\gamma,K_\infty)}.
 \label{eq:q4-sign-linking-parity}
\end{equation}
Consequently, the full inverse image of \(K_\gamma\) in
\(S^3\setminus T(2,4)\) has two components when
\(m+\operatorname{Lk}(K_\gamma,K_\infty)\) is even and one component when it
is odd.
\end{corollary}

\begin{proof}
Every \(u\)-turn is the nontrivial rotation at an order-two vertex, hence a
conjugate of \(J\).  Every \(v^{j_k}\)-turn is, in the moving vertex frame, a
conjugate of \(R^{j_k}\) or \(R^{-j_k}\), according to the clockwise
orientation convention.  Because \(\chi\) is conjugacy invariant and
\(\chi(R^{\pm j_k})=(-1)^{j_k}\), multiplying the turn contributions gives
\(\chi(\gamma)=\prod_k(-1)^{1+j_k}\).  This argument uses the classical
geometric code-to-group calibration; it does not identify the code with a
literal Paper~I history word.
On the other hand, Theorem~\ref{thm:q4-cusp-linking-code} gives
\[
 \operatorname{Lk}(K_\gamma,K_\infty)
 =\sum_k j_k-2m\equiv\sum_k j_k\pmod2.
\]
The last assertion follows from
Proposition~\ref{prop:q4-sign-controlled-knot-lift}.
\end{proof}

\begin{theorem}[The zero dessin as the classical coding spine]
\label{thm:q4-zero-dessin-coding-spine}
Match the fundamental triangle used to normalize \(\beta\) with the triangle
used in the classical \((2,4,\infty)\) cutting construction.  Then:
\begin{enumerate}
\item the Dehornoy--Pinsky graph \(G[2,4,\infty]\) is exactly the AEG zero
dessin,
\begin{equation}
 G[2,4,\infty]
 =\beta^{-1}([0,1])
 =Z(a_4),
 \label{eq:q4-coding-graph-zero-graph}
\end{equation}
as an embedded \(\Gamma_4\)-graph in \(\HH\);
\item this graph is the barycentric subdivision of the four-regular tree dual
to the adapted tessellation by ideal quadrilaterals; suppressing its bivalent
order-two vertices recovers that dual tree;
\item the two endpoints of a hyperbolic axis determine the unique
Dehornoy--Pinsky admissible bi-infinite path when \(r=\infty\).  For an axis
stabilized by \(\gamma\), this path is periodic modulo \(\gamma\); after
suppressing the bivalent vertices it is a no-backtracking path in the dual zero
tree, and its turns give the cyclic code
\eqref{eq:q4-reduced-cutting-code};
\item replacing the admissible graph path by the classical roundabout path and
then thickening the allowed visual-direction intervals produces the template.
The unit-tangent knot \(K_\gamma\) is isotopic to the corresponding template
orbit, simultaneously for every finite collection of periodic orbits.
\end{enumerate}
\end{theorem}

\begin{proof}
The graph \(G[2,4,\infty]\) is defined as the
\(\Gamma_4\)-orbit of the geodesic segment joining the order-two and order-four
vertices.  The proof of Theorem~\ref{thm:q4-hecke-zero-network} identifies
\(\beta^{-1}([0,1])\) with the orbit of the same segment, which proves
\eqref{eq:q4-coding-graph-zero-graph}.  Suppressing the order-two midpoint of
each pair of adjacent ideal quadrilaterals gives their ordinary dual tree.

The remaining assertions are the classical coding and template theorems.
For \(r=\infty\) the coding graph is a tree and the spectacles choice is unique.
It assigns to the two ideal endpoints of every geodesic an admissible graph
path, including the boundary cases for which a naive crossed-tile itinerary
would be ambiguous.  Periodicity of the geodesic is equivalent to periodicity
of the reduced path code.  The classical roundabout replacement records the
entry and exit data inside each tile; retaining the corresponding interval of
unit tangent directions produces the ribbons.  The tearing construction then
gives the asserted simultaneous isotopy
\cite{Dehornoy2015GeodesicFlow,DehornoyPinsky2018Coding}.
\end{proof}

\begin{warning}[The spine is not the template]
\label{warn:q4-spine-not-template}
The hyperbolic axis is not literally contained in \(Z(a_4)\).  Its tile
itinerary is represented on the dual zero tree.  Moreover, the bare tree does
not remember the visual direction needed to select the chord and the
over--under placement of ribbons in the unit tangent bundle.  Thus
Theorem~\ref{thm:q4-zero-dessin-coding-spine} does not identify an AEG zero tube
with the classical template; it identifies the AEG zero dessin with the
one-dimensional coding spine from which the directed template is constructed.
\end{warning}

\begin{proposition}[The operator-quotient history-to-knot map]
\label{prop:q4-operator-history-knot-map}
Let \(\operatorname{Hist}_{4}^{\mathrm{ph}}\) be the Paper~I words in
\(\mathsf t_+,\mathsf t_-\), and \(\mathsf j\) whose projective operator is
primitive and hyperbolic.  Then
\begin{equation}
 \mathcal K_{\mathrm{op}}:
 \operatorname{Hist}_{4}^{\mathrm{ph}}
 \longrightarrow
 \{\text{oriented primitive geodesic-flow periodic orbits in }M_4\},
 \qquad
 h\longmapsto K_{\rho(h)},
 \label{eq:q4-operator-history-knot-map}
\end{equation}
is surjective and has the following invariance properties.
\begin{enumerate}
\item A cyclic shift of the chronological word does not change the oriented
knot: its operator is conjugate to \(\rho(h)\).
\item Reversing the word and inverting every letter replaces the orbit by the
time-reversed orbit \(K_{\rho(h)^{-1}}\), which in general is not merely the
same embedded circle with its orientation reversed.
\item The map factors through the conjugacy class of the projective operator.
In particular, inserting Paper~I's neutral eight-letter word \(\omega_4\), or
making any other marked-history change that preserves that conjugacy class,
cannot be detected by the unframed knot.
\end{enumerate}
No ordinary real-input admissibility is asserted by
\eqref{eq:q4-operator-history-knot-map}.
\end{proposition}

\begin{proof}
Write a chronological word as \(h=(g_1,\ldots,g_n)\), with matrix
representatives \(M_i\).  Paper~I's convention gives
\(\rho(h)=[M_n\cdots M_1]\).  The cyclic shift
\((g_2,\ldots,g_n,g_1)\) has operator
\[
 [M_1M_n\cdots M_2]
 =[M_1]\rho(h)[M_1]^{-1},
\]
so Proposition~\ref{prop:q4-hyperbolic-class-knot} proves the first claim.
The path inverse evaluates to the inverse projective operator, proving the
second.  The third is then immediate, and
\(\rho(\omega_4)=1\) by Paper~I.  Finally, every element of \(\Gamma_4\) is
represented by this history alphabet; hence every primitive hyperbolic class,
and therefore every knot in the target, has a preimage.
\end{proof}

\begin{problem}[Enhancing operator histories beyond the knot quotient]
\label{prob:q4-history-geodesic-naturality}
Proposition~\ref{prop:q4-operator-history-knot-map} settles the unadorned
operator-level map but also exhibits its large kernel.  Construct a declared
marked cyclic equivalence and a canonical enhancement
\[
 \operatorname{Hist}_{4}^{\mathrm{ph}}/\mathord{\sim}_{\mathrm{cyc}}
 \longrightarrow
 \{\text{classical cyclic cutting codes with AEG history decoration}\}
\]
that specifies ordinary-arithmetic admissibility, base edge, orientation, and
the residual data carried by histories with conjugate operators.  It must prove
compatibility with the cutting construction of
Theorem~\ref{thm:q4-zero-dessin-coding-spine} and characterize the marked and
syntactic residual structure inside the conjugacy-class fibers of
\(\mathcal K_{\mathrm{op}}\), or the still larger fibers after forgetting to
ambient isotopy type.  Without such an enhancement, the classical code
is an intrinsic code of the Hecke operator class, not of the full marked AEG
history.
\end{problem}

\begin{remark}[Current novelty boundary]
Theorems~\ref{thm:q4-sign-cover-torus-link} and
\ref{thm:q4-cusp-linking-code} give a rigorous route
\[
 \begin{gathered}
 \text{Hecke conjugacy and cutting data}
 \longrightarrow \text{closed geodesics}\\
 \longrightarrow \text{knots and linking numbers}.
 \end{gathered}
\]
The character \(\chi\) additionally controls how the knots split in the natural
\(S^3\) cover.  The classical geodesic-flow and linking parts are not new knot
invariants.  The AEG-specific frontier is to enhance
Proposition~\ref{prop:q4-operator-history-knot-map} as required by
Problem~\ref{prob:q4-history-geodesic-naturality}, and then determine whether the
assignment, cone, sign, divisor, or history data define an isotopy invariant not
already determined by the classical Hecke code or its known Rademacher-type
functions.  Nothing in this section asserts such a separation theorem.
\end{remark}
